% Options for packages loaded elsewhere
\PassOptionsToPackage{unicode}{hyperref}
\PassOptionsToPackage{hyphens}{url}
\documentclass[
  ignorenonframetext,
]{beamer}
\newif\ifbibliography
\usepackage{pgfpages}
\setbeamertemplate{caption}[numbered]
\setbeamertemplate{caption label separator}{: }
\setbeamercolor{caption name}{fg=normal text.fg}
\beamertemplatenavigationsymbolsempty
% remove section numbering
\setbeamertemplate{part page}{
  \centering
  \begin{beamercolorbox}[sep=16pt,center]{part title}
    \usebeamerfont{part title}\insertpart\par
  \end{beamercolorbox}
}
\setbeamertemplate{section page}{
  \centering
  \begin{beamercolorbox}[sep=12pt,center]{section title}
    \usebeamerfont{section title}\insertsection\par
  \end{beamercolorbox}
}
\setbeamertemplate{subsection page}{
  \centering
  \begin{beamercolorbox}[sep=8pt,center]{subsection title}
    \usebeamerfont{subsection title}\insertsubsection\par
  \end{beamercolorbox}
}
% Prevent slide breaks in the middle of a paragraph
\widowpenalties 1 10000
\raggedbottom
\AtBeginPart{
  \frame{\partpage}
}
\AtBeginSection{
  \ifbibliography
  \else
    \frame{\sectionpage}
  \fi
}
\AtBeginSubsection{
  \frame{\subsectionpage}
}
\usepackage{iftex}
\ifPDFTeX
  \usepackage[T1]{fontenc}
  \usepackage[utf8]{inputenc}
  \usepackage{textcomp} % provide euro and other symbols
\else % if luatex or xetex
  \usepackage{unicode-math} % this also loads fontspec
  \defaultfontfeatures{Scale=MatchLowercase}
  \defaultfontfeatures[\rmfamily]{Ligatures=TeX,Scale=1}
\fi
\usepackage{lmodern}
\ifPDFTeX\else
  % xetex/luatex font selection
\fi
% Use upquote if available, for straight quotes in verbatim environments
\IfFileExists{upquote.sty}{\usepackage{upquote}}{}
\IfFileExists{microtype.sty}{% use microtype if available
  \usepackage[]{microtype}
  \UseMicrotypeSet[protrusion]{basicmath} % disable protrusion for tt fonts
}{}
\makeatletter
\@ifundefined{KOMAClassName}{% if non-KOMA class
  \IfFileExists{parskip.sty}{%
    \usepackage{parskip}
  }{% else
    \setlength{\parindent}{0pt}
    \setlength{\parskip}{6pt plus 2pt minus 1pt}}
}{% if KOMA class
  \KOMAoptions{parskip=half}}
\makeatother
\usepackage{longtable,booktabs,array}
\usepackage{caption}
\captionsetup[table]{skip=6pt}
\newcounter{none} % for unnumbered tables
\usepackage{calc} % for calculating minipage widths
\usepackage{caption}
% Make caption package work with longtable
\makeatletter
\def\fnum@table{\tablename~\thetable}
\makeatother
\setlength{\emergencystretch}{3em} % prevent overfull lines
\providecommand{\tightlist}{%
  \setlength{\itemsep}{0pt}\setlength{\parskip}{0pt}}
% Auto-generated by infrastructure/rendering/_pdf_latex_helpers.py::extract_math_font_preamble.
% Minimal header passed to Pandoc via -H so Beamer slides pick up the same math font
% as the combined-PDF path. Adjust the manuscript's preamble.md to change the font globally.
\usepackage{unicode-math}
\setmathfont{latinmodern-math.otf}

% Manuscript macro/package fallbacks (newcommand -> providecommand).
% Core mathematics
\usepackage{amsmath}
\usepackage{amssymb}
% Document layout
% Tables
\usepackage{booktabs}
% Algorithm typesetting (for pseudocode in §2 Methodology)
% Code listings
\usepackage{listings}
% Typography and formatting
% Cross-references and citations
% ── Unicode-capable mono font for code listings ──────────────────────
% JuliaMono (TeX Live 2026) covers the full math/Greek glyph set used in
% scientific code blocks; the default lmmono lacks \alpha/\mu/\partial/\nabla.
%
% The hardcoded TeX Live 2026 path below is portable to most macOS / Linux
% TeX Live installs that placed JuliaMono under the default texmf-dist path.
% If your TeX Live version differs (e.g., 2025, 2027) or your distribution
% installs fonts elsewhere, comment out the explicit Path = line — fontspec
% will then search the system font path. If JuliaMono is not installed at
% all, comment out the whole block; xelatex falls back to lmmono with
% reduced Unicode coverage but the document still compiles.
% Math font for unicode-math: Latin Modern Math (TeX Live) has full BMP
% coverage including U+2223 (\mid), U+226A/226B (\ll/\gg), and the Greek/
% blackboard letters used in equations. Without an explicit \setmathfont,
% unicode-math falls back to lmroman text font which lacks several glyphs
% and emits "Missing character" warnings on every \mid in math mode.

% Slide layout defaults for warning-clean scientific decks.
\usepackage{etoolbox}
\IfFileExists{xurl.sty}{\usepackage{xurl}}{}
\IfFileExists{seqsplit.sty}{\usepackage{seqsplit}}{\newcommand{\seqsplit}[1]{#1}}
\protected\def\breakseq#1{\seqsplit{#1}}
\protected\def\breaktt#1{\begingroup\ttfamily\seqsplit{#1}\endgroup}
\setlength{\emergencystretch}{6em}
\tolerance=5000
\hbadness=10000
\hfuzz=1pt
\setlength{\tabcolsep}{2pt}
\AtBeginEnvironment{longtable}{\tiny\renewcommand{\arraystretch}{0.86}\setlength{\tabcolsep}{1pt}}
\AtBeginEnvironment{tabular}{\tiny\renewcommand{\arraystretch}{0.86}\setlength{\tabcolsep}{1pt}}
\AtBeginEnvironment{equation}{\tiny}
\AtBeginEnvironment{equation*}{\tiny}
\AtBeginEnvironment{align}{\tiny}
\AtBeginEnvironment{align*}{\tiny}
\AtBeginEnvironment{itemize}{\footnotesize}
\AtBeginEnvironment{enumerate}{\footnotesize}
\AtBeginEnvironment{description}{\footnotesize}
\setbeamerfont{caption}{size=\tiny}
\setbeamerfont{caption name}{size=\tiny}
\setbeamerfont{normal text}{size=\small}
\setbeamerfont{frametitle}{size=\small}
\setbeamerfont{section title}{size=\footnotesize}
\setbeamerfont{subsection title}{size=\footnotesize}
\setbeamertemplate{section page}{%
  \centering
  \begin{beamercolorbox}[sep=12pt,center,wd=\paperwidth]{section title}
    \parbox{0.86\paperwidth}{\centering\usebeamerfont{section title}\insertsection\par}
  \end{beamercolorbox}
}
\setbeamertemplate{subsection page}{%
  \centering
  \begin{beamercolorbox}[sep=8pt,center,wd=\paperwidth]{subsection title}
    \parbox{0.86\paperwidth}{\centering\usebeamerfont{subsection title}\insertsubsection\par}
  \end{beamercolorbox}
}
\setlength{\abovecaptionskip}{2pt}
\setlength{\belowcaptionskip}{0pt}

% Natbib and cross-reference fallbacks — slides don't load natbib
% or cleveref, but combined-PDF manuscript prose may emit these
% commands. The fallback renders citations as a bracketed key list
% and cross-references as detokenized labels so slides don't fail on
% undefined control sequences or raw underscores. \providecommand is
% a no-op if packages load later.
\providecommand{\citep}[1]{[#1]}
\providecommand{\citet}[1]{#1}
\providecommand{\citealp}[1]{#1}
\providecommand{\citeauthor}[1]{#1}
\providecommand{\citeyear}[1]{#1}
\providecommand{\cref}[1]{\texttt{\detokenize{#1}}}
\providecommand{\Cref}[1]{\texttt{\detokenize{#1}}}
\providecommand{\crefrange}[2]{\texttt{\detokenize{#1}}--\texttt{\detokenize{#2}}}
\providecommand{\Crefrange}[2]{\texttt{\detokenize{#1}}--\texttt{\detokenize{#2}}}
\providecommand{\paragraph}[1]{\textbf{#1}\ }

\newtheorem{proposition}{Proposition}
\newtheorem{hypothesis}{Hypothesis}
\newtheorem{remark}[theorem]{Remark}
\newtheorem{axiom}{Axiom}
\newtheorem{property}{Property}
\usepackage{bookmark}
\IfFileExists{xurl.sty}{\usepackage{xurl}}{} % add URL line breaks if available
\urlstyle{same}
% fallback for those not using the hyperref driver hyperxmp:
\makeatletter
\@ifundefined{xmpquote}{\newcommand{\xmpquote}[1]{#1}}{}
\makeatother
\hypersetup{
  hidelinks,
  pdfcreator={LaTeX via pandoc}}

\author{\texorpdfstring{}{}}
\date{}

\begin{document}

\section{Reproducible Operations: NixOS in Depth}\label{sec:nixos}

\begin{frame}[fragile,allowframebreaks]{The model: declarative
configuration, generations, rollback}
\protect\phantomsection\label{the-model-declarative-configuration-generations-rollback}
NixOS expresses an entire system --- packages, services, users, kernel
settings --- as a declarative configuration evaluated by the Nix package
manager, and activates changes as numbered \emph{generations}
({community} 2026e). Activation builds a new system closure, atomically
switches to it, and leaves prior generations bootable, so
\texttt{nixos-rebuild} can roll the system configuration back to an
earlier point ({community} 2026f). Nix builds run in isolated filesystem
and process namespaces, making controlled build inputs a starting point
for reproducible construction, though the reproducibility project
documents remaining sources of nondeterminism (N. Project 2026g, 2026j).

\end{frame}

\begin{frame}[fragile,allowframebreaks]{The model: declarative
configuration, generations, rollback}
The package manager itself has been evolving on a cadence that matters
for security analysis. Nix 2.34 (released early 2026) tightened sandbox
behavior --- blocking extended-attribute propagation inside builds ---
and extended build tracing, and Nix 2.35.0 (June 22, 2026) continues
that line; the 2.35 release notes also fold secrets handling into the
manual as a first-class chapter, recognizing that credential delivery
into builds and services is a security surface rather than an
afterthought (N. Project 2026f). Version discipline is part of the
update-operations property: an installation's security posture is a
function of which Nix release it actually runs, and the advisory record
below shows the Nix releases themselves requiring prompt updates.

\end{frame}

\begin{frame}[fragile,allowframebreaks]{The model: declarative
configuration, generations, rollback}
In this review's threat model, NixOS's value is not confinement. It is
the ability to make intended system state explicit, auditable, and
replaceable --- to construct, update, audit, and replace tightly scoped
environments that are separately isolated. A deployment whose
configuration is a reviewed, versioned artifact can answer ``what
exactly is running here?'' in a way an imperatively mutated machine
cannot, which is the foundation the authority architecture in
sec.~10 builds on when an agent proposes
configuration changes.
\end{frame}

\begin{frame}[fragile,allowframebreaks]{Rollback is not recovery}
\protect\phantomsection\label{rollback-is-not-recovery}
Generation rollback switches \emph{system configuration} generations; it
does not restore mutable state ({community} 2026f). The community
discussion of rolling back data as well as configuration documents the
mismatch that arises when a service's database has migrated but its
software is rolled back ({community} 2026j). Three consequences follow
for the persistence-and-recovery property.

First, \textbf{rollback is not a malware-remediation plan}. A
compromised agent's files in \texttt{/home}, stolen credentials,
exfiltrated data, and accepted remote changes survive a generation
switch; the state a rebuild can restore is the state Nix owns. Second,
\textbf{older retained generations are potentially vulnerable software,
not trusted recovery images}: they are pinned to the package versions of
their era, including any advisories published since --- a concern the
2026 Nix advisories make concrete, since a generation built with an
\end{frame}

\begin{frame}[fragile,allowframebreaks]{Rollback is not recovery}
unpatched Nix release carries the advisory's exposure until the manager
itself is updated. Third, \textbf{data recovery is a separate
discipline} --- independent backups and rehearsed restoration
(sec.~12) --- that declarative configuration does not
substitute for. The review's recovery scenario in
sec.~16 therefore requires distinguishing deletion of
an environment from revocation of any credentials it could have reached.
\end{frame}

\begin{frame}[fragile,allowframebreaks]{Reproducible deployment versus
verified reproducible builds}
\protect\phantomsection\label{reproducible-deployment-versus-verified-reproducible-builds}
The most misused word in this space is ``reproducible,'' and NixOS is
where the confusion does real security harm. Two distinct claims must be
kept apart (N. Project 2026j):

\end{frame}

\begin{frame}[fragile,allowframebreaks]{Reproducible deployment versus
verified reproducible builds}
\begin{itemize}
\tightlist
\item
  \textbf{Reproducible deployment}: a system can be recreated from its
  declarative configuration --- the same expression yields a coherent,
  working system. NixOS provides this well ({community} 2026e).
\item
  \textbf{Verified reproducible builds}: every package has been
  independently rebuilt and compared bit-for-bit against the binary
  actually deployed, proving the artifact corresponds to its claimed
  inputs.
\end{itemize}

\end{frame}

\begin{frame}[fragile,allowframebreaks]{Reproducible deployment versus
verified reproducible builds}
The first does not imply the second. The Nix sandbox's dependency
specification and filesystem/process isolation alone do not guarantee
reproducible outputs (N. Project 2026g, 2026j), and the reproducibility
project exists precisely because the gap is being closed package by
package, not closed by declaration. The 2026 measurements quantify how
far the gap has closed and where it remains: the NixOS ISO closure
(channel \texttt{nixos.iso\_gnome}) reaches 95.18\% reproducibility
across its full paths (N. Project 2026h), while a build-closure analysis
reports 99.48\% of packages rebuilding bit-for-bit (N. Project 2026h); a
large-scale rebuild study of Nixpkgs documents the remaining
nondeterminism classes and their distribution across the package set
(Malka et al. 2025). These are measured rates for specific artifacts,
not a general guarantee --- the same discipline that keeps this review
from scoring platforms applies to interpreting them.
\end{frame}

\begin{frame}[fragile,allowframebreaks]{Reproducible deployment versus
verified reproducible builds}

Neither form of reproducibility, finally, establishes that the source,
configuration, or upstream dependency is \emph{safe} --- reproducibility
answers ``is this the artifact that source produces?'', not ``should
anyone run this?''. This mirrors the review's broader separations:
authenticity from provenance, rollback from recovery, and containment
from authority (sec.~4).
\end{frame}

\begin{frame}[fragile,allowframebreaks]{Six misconceptions that matter
against offensive agents}
\protect\phantomsection\label{six-misconceptions-that-matter-against-offensive-agents}
Because Nix's vocabulary sounds stronger than its documented guarantees,
the source assessment's misconception list is worth keeping as a
standing table. tbl.~\ref{tbl:nixos_misconceptions} states each with the
documented fact and its consequence --- updated where the 2026 record
changed the facts.

\end{frame}

\begin{frame}[fragile,allowframebreaks]{Six misconceptions that matter
against offensive agents}
\begin{longtable}[]{@{}
  >{\raggedright\arraybackslash}p{(\linewidth - 4\tabcolsep) * \real{0.3333}}
  >{\raggedright\arraybackslash}p{(\linewidth - 4\tabcolsep) * \real{0.3333}}
  >{\raggedright\arraybackslash}p{(\linewidth - 4\tabcolsep) * \real{0.3333}}@{}}
\caption{Six NixOS misconceptions, the documented facts that correct
them, and the consequences against offensive
agents.}\label{tbl:nixos_misconceptions}\tabularnewline
\toprule\noalign{}
\begin{minipage}[b]{\linewidth}\raggedright
Misconception
\end{minipage} & \begin{minipage}[b]{\linewidth}\raggedright
Documented fact
\end{minipage} & \begin{minipage}[b]{\linewidth}\raggedright
Consequence
\end{minipage} \\
\midrule\noalign{}
\endfirsthead
\toprule\noalign{}
\begin{minipage}[b]{\linewidth}\raggedright
Misconception
\end{minipage} & \begin{minipage}[b]{\linewidth}\raggedright
Documented fact
\end{minipage} & \begin{minipage}[b]{\linewidth}\raggedright
Consequence
\end{minipage} \\
\midrule\noalign{}
\endhead
``The Nix sandbox confines everything I run.'' & The documented sandbox
concerns builds, including specified exceptions such as network access
for fixed-output derivations (N. Project 2026g). & The sandbox is a
build-time mechanism, not a runtime security boundary around an AI agent
or programs installed through Nix. \\
``A signed cache artifact is proven to match reviewed source.'' & Cache
signatures distinguish trusting a signer from proving an output was
produced by the expected derivation ({community} 2026m). & Authenticity
is valuable, but a trusted supplier can still distribute a bad artifact;
signatures do not establish provenance --- which is why the provenance
tooling below matters. \\
``Putting secrets into declarative configuration is harmless.'' & The
Nix manual treats secret management as a first-class concern: the store
is world-readable and secrets can reach external caches, so secrets
belong in runtime-protected stores --- now with dedicated manual
guidance in the 2.35 documentation (N. Project 2026k, 2026f). &
Configuration cleanliness must not become credential exposure; an agent
with store access can read declaratively embedded secrets. \\
``NixOS already has a fully integrated mandatory-access-control
baseline.'' & The wiki documents incomplete SELinux integration and, as
of April 2026, incomplete AppArmor integration ({community} 2026k). &
Confinement is possible but not presumed; do not assume a comprehensive
default policy exists. \\
``Secure Boot is a settled, integrated default.'' & Lanzaboote remains
outside nixpkgs --- its packaging PR is still open ({community} 2026b)
--- and its repository documents that key management is outside the
project's full scope; measured-boot (TPM PCR) support landed in April
2026, and the bootspec removal in nixpkgs broke v1.0.0 until v1.1.0
fixed the integration ({community} 2026a, 2026c). & A carefully
configured installation and a stock default are different claims;
boot-integrity stance is \texttt{partial} pending upstream integration
--- and the version pinning is load-bearing. \\
``The hardened profile is a safe universal upgrade.'' & The hardened
profiles were \textbf{removed} from nixpkgs on March 22, 2026 (PR
nixpkgs\#501199), with the 26.05 release notes listing the removal as a
backward incompatibility; the maintainer discussion that preceded it
argued the profile lacked a coherent baseline and could undermine
browser sandboxing, and the NixOS Hardening wiki now carries the
composable guidance instead ({community} 2026h, 2026d, 2026g). &
Hardening menus are not defaults; the ecosystem itself concluded that an
unmaintained blanket profile is worse than documented per-concern
guidance --- the defaults-over-menus forecast in
sec.~15. \\
\bottomrule\noalign{}
\end{longtable}

\end{frame}

\begin{frame}[fragile,allowframebreaks]{Six misconceptions that matter
against offensive agents}
For agent-bearing deployments the first and third rows bind hardest: an
autonomous agent executing on a NixOS host is \emph{not} sandboxed by
Nix, and any secret that reached declarative configuration is readable
by whatever that agent's processes can read.
\end{frame}

\begin{frame}[fragile,allowframebreaks]{The package manager is part of
the trusted computing base}
\protect\phantomsection\label{the-package-manager-is-part-of-the-trusted-computing-base}
The 2026 advisory chain makes the trusted-computing-base argument with
unusual clarity, because the advisories landed in sequence against the
same component. On April 7, 2026, Nix published a security advisory
describing a symlink-following flaw during fixed-output derivation
registration that could allow users permitted to submit builds to gain
root in multi-user installations, including with sandboxed Linux builds
enabled; the advisory lists fixes across version branches, including
2.34.5 and 2.33.4 (N. Project 2026b). Later in the year two further
advisories followed: one in the NAR archive parser, where unbounded
recursion --- now bounded --- allowed resource exhaustion during
store-path ingestion at high severity (CVSS 7.5) (N. Project 2026d); and
one enabling path traversal on \texttt{-\/-unpack} extraction, an
\end{frame}

\begin{frame}[fragile,allowframebreaks]{The package manager is part of
the trusted computing base}
attacker-controlled path escape out of the intended destination (CVSS
4.3) (N. Project 2026c). A third issue, a TOCTOU race in recursive-Nix
builds, required the 2.35.0 fix to close (N. Project 2026a, 2026f). Each
advisory is fixed in current releases; the structural observation is
what matters.

The finding is not evidence that Nix is uniquely unsafe, nor that the
flaws remain unpatched. It is evidence for a structural claim:
\textbf{the build service is part of the trusted computing base}. A
component with root authority over the machine, whose input includes
build submissions and attacker-named artifacts (NAR archives, tarballs,
path names), is a privileged parser of attacker-influenced content ---
the same shape as the QSB-118 dom0 injection in sec.~5,
and the same lesson: prompt patching of the management layer is not
optional hygiene but the update-operations property of
sec.~4 doing real security work. The chain
\end{frame}

\begin{frame}[fragile,allowframebreaks]{The package manager is part of
the trusted computing base}
also illustrates the second-order risk of pinned environments: a
deployment frozen on an older Nix release accumulates advisory exposure
in the very tool that delivers its other updates.

The deployment consequence follows directly from the threat model in
sec.~3: an agent allowed to submit
attacker-controlled builds should not share a trust domain with the
owner's most valuable credentials merely because the builds use Nix. In
the trust-domain architecture of sec.~10, build
submission belongs with agent execution (disposable, scoped), not with
administration or the credential service --- and a multi-user Nix
installation hosting untrusted build submitters inherits a
root-equivalent trust requirement that its operators must patch and
monitor (N. Project 2026b, 2026d). The Nixpkgs security tracker (N.
Project 2026i) coordinates vulnerability matching and mitigation, which
supports the update-operations stance but does not by itself deliver
\end{frame}

\begin{frame}[fragile,allowframebreaks]{The package manager is part of
the trusted computing base}
patch latency.
\end{frame}

\begin{frame}[fragile,allowframebreaks]{Provenance: from signatures to
verification}
\protect\phantomsection\label{provenance-from-signatures-to-verification}
The supply-chain-trust property asks who may introduce or approve
executable code, and 2026 hardened the tooling that answers it. Nixpkgs
now requires \breaktt{meta.sourceProvenance} on packages --- an explicit,
machine-readable declaration of whether a build artifact derives from
source, from an upstream binary blob, or from an unfree distribution ---
turning provenance from documentation into an enforced field of every
package definition ({community} 2026i). Verification tooling closes the
loop on the binary side: sbomnix generates SBOMs (CycloneDX/SPDX) for a
Nix closure and can attach SLSA provenance attestations ((TII) 2026);
Trustix compares independently rebuilt derivations at scale, exposing
bit-for-bit divergence between builders ({community} 2026l); and
Determinate's \breaktt{nix\ provenance\ show}/\texttt{verify} commands
\end{frame}

\begin{frame}[fragile,allowframebreaks]{Provenance: from signatures to
verification}
bring attestation checking into the package-manager CLI itself. None of
this converts a signature into proof of benignity --- the misconception
table's second row still stands --- but together they make the
provenance half of the question answerable by tooling rather than trust,
which is exactly the direction the supply-chain-trust property demands
for agent-bearing deployments: an auditor can now ask ``what does this
closure claim, and what was independently rebuilt?'' as data.
\end{frame}

\begin{frame}[allowframebreaks]{Update operations: the support-clock
discipline}
\protect\phantomsection\label{update-operations-the-support-clock-discipline}
NixOS 26.05's announcement specifies seven months of bug fixes and
security updates, ending December 31, 2026 (N. Project 2026e). A stable
configuration revision is a point-in-time choice, not a permanently safe
one: after support ends, an unchanged deployment accumulates unpatched
vulnerabilities while its configuration remains byte-identical. The
update-operations property therefore binds declarative systems as hard
as any other: establish an explicit channel-update and rebuild cadence,
treat held-back generations as aging software, and fold the tracker into
the operator's patch process (N. Project 2026i, 2026e). This is where
Nix's strengths compound --- because the system state is declarative, an
update is a \emph{reviewable configuration change} rather than an opaque
mutation, and a failed rebuild leaves the prior generation bootable
\end{frame}

\begin{frame}[allowframebreaks]{Update operations: the support-clock
discipline}
({community} 2026f). The 2026 advisory chain gives the discipline a
concrete rhythm: each Nix release that closes an advisory is also a
reminder that the update path itself is the artifact being protected.
\end{frame}

\begin{frame}[allowframebreaks]{Composition with Qubes}
\protect\phantomsection\label{composition-with-qubes}
Nix and Qubes are not competing answers: one principally organizes
runtime trust domains, the other organizes how software environments are
constructed. The natural composition uses Nix to construct the contents
of compartmentalized qubes --- declarative, reviewable, replaceable
environments inside a containment architecture. The Qubes NixOS-template
issue documents community work toward exactly this (Q. O. Project
2026a), but the integration is not a product: Qubes warns that community
templates do not receive updates from the Qubes project itself (Q. O.
Project 2026b), so a hand-built Qubes-plus-NixOS stack inherits a
maintenance obligation for template security updates, boot support, and
desktop integration that its owner must explicitly own ({community}
2026k, 2026b). Combining two attractive designs does not automatically
improve practical security; the composition must be operated, patched,
\end{frame}

\begin{frame}[allowframebreaks]{Composition with Qubes}
and reviewed as its own deployment.
\end{frame}

\begin{frame}[fragile,allowframebreaks]{Assessment}
\protect\phantomsection\label{assessment}
NixOS is especially attractive for a capable operator building
auditable, replaceable development or server environments. On the
candidate--property matrix its update-operations and supply-chain-trust
stances are \texttt{strong} --- declarative state, channel discipline,
the security tracker, and now enforced provenance metadata and
verification tooling are documented and coherent ({community} 2026i;
(TII) 2026) --- while its application-confinement stance is
\texttt{weak} pending integrated MAC and boot-integrity baselines
({community} 2026k, 2026a). It is not the default recommendation over
Qubes for containing a hostile desktop workload, and it is not
sufficient by itself to make a powerful autonomous coding agent safe:
the sandbox that defines its build discipline is build-time, its secrets
model demands deliberate handling --- now with first-class manual
\end{frame}

\begin{frame}[fragile,allowframebreaks]{Assessment}
guidance (N. Project 2026f) --- and its build service is root-authority
code with a 2026 advisory chain that must be treated as trusted base (N.
Project 2026g, 2026k, 2026b, 2026d, 2026c). The strongest arrangement
--- and the one this review's architecture chapters target --- uses
Nix's declarative construction \emph{inside} a separately enforced,
least-privilege execution architecture, where generation discipline
supplies replaceability and the containment layer supplies what Nix
deliberately does not.
\end{frame}

\end{document}
